\section{MIDI}\label{midi}

The synthesizer has 2 MIDI interfaces: a standard serial interface using
DIN connectors, and a MIDI USB interface for use with a computer. Both
provide an input/output pair.

\subsection{MIDI Channel}\label{midi-channel}

The MIDI channel (1--16) is selected in the settings menu. The synth
listens to incoming messages on this channel and sends sequencer notes
to the MIDI output on the same channel.

MIDI input can be disabled per voice track by setting its MIDI channel
to `0' in the global settings, or completely disabled by setting the
global MIDI channel to `0'.

\subsection{MIDI Sync}\label{midi-sync}

The LXR can be synced to external gear using MIDI clock signals. Set the
BPM parameter in the settings menu to \texttt{0} to enable slave mode.
The LXR will then sync its tempo to incoming MIDI clocks, and the Play
button has no function --- the synth listens for MIDI start/stop signals
instead.

When not in MIDI slave mode, the LXR sends its internal clock as MIDI
clock messages on both the USB and the DIN MIDI output. This way the LXR
can act as a MIDI master and sync other devices. MIDI start/stop
messages are sent when the sequencer starts or stops.

\subsection{Record Incoming MIDI}\label{record-incoming-midi}

You can trigger all voices with external MIDI gear (drum pads, keyboard,
etc.). If record mode is active, the played MIDI notes are directly
recorded into the corresponding patterns. Note recording can be
quantised to a grid.

{[}\textbf{Brendan: To Review}{]} \textbf{Global MIDI Note Trigger:} If
a MIDI note is received on the designated global MIDI channel, the LXR
will trigger whichever track corresponds to the active track selected on
the front panel, unless that track has a specific note override
configured. This allows you to play the selected instrument
chromatically from an external keyboard.

The default MIDI note numbers for the 7 tracks are:

\begin{voicemiditbl}
1 & 36 \\
2 & 37 \\
3 & 38 \\
4 & 39 \\
5 & 40 \\
6 & 41 \\
7 & 42 \\
\end{voicemiditbl}

\textbf{MIDI roll trigger:} when the MIDI Roll Offset setting
(\param{rol} in the Global Settings menu) is on, a note that many
semitones above a voice's trigger note plays a held roll on that voice
instead of a single hit. See \emph{MIDI Roll Offset} in the Global
Settings Menu section. A note-on with velocity 0 is treated as a
note-off.

\subsection{Transmitting Sequencer Data to the MIDI
Output}\label{transmitting-sequencer-data-to-the-midi-output}

The sequencer currently transmits note-on trigger messages to the MIDI
output. The ability to send complete sequencer tracks with pitch
information on separate MIDI channels is planned for a future release.

\subsection{Transmitting Parameter
Changes}\label{transmitting-parameter-changes}

Adjustments to voice parameters made in the menu are transmitted on the
Global MIDI channel when the \texttt{txf} global setting includes
\textbf{C}.

\subsection{Program / Bank Changes}\label{program-bank-changes}

Program Change messages on the global channel select the next pattern
(0-7) for all tracks. Program Change messages on voice channels select
the pattern for that track independently.

Bank Select (CC0) and Mod Wheel (CC1) are supported. CC0 triggers
bank/kit changes, and CC1 controls global morph on the global MIDI
channel, or voice morph on individual voice MIDI channels.

\subsection{MIDI CC List}\label{midi-cc-list}

Parameters can be remote-controlled with MIDI CC messages.
{[}\textbf{Brendan: To Review}{]} The LXR now supports independent MIDI
CC assignments per voice via a dedicated global channel, as well as the
legacy per-channel mapping.

\subsubsection{Channel MIDI CC Table}\label{channel-midi-cc-table}

If a CC is received on a voice's specific MIDI channel, it controls that
voice according to the table below.

{\def\LTcaptype{none} % do not increment counter
\begin{longtable}[]{@{}
  >{\raggedleft\arraybackslash}p{(\linewidth - 16\tabcolsep) * \real{0.1429}}
  >{\raggedright\arraybackslash}p{(\linewidth - 16\tabcolsep) * \real{0.1071}}
  >{\raggedright\arraybackslash}p{(\linewidth - 16\tabcolsep) * \real{0.1071}}
  >{\raggedright\arraybackslash}p{(\linewidth - 16\tabcolsep) * \real{0.1071}}
  >{\raggedright\arraybackslash}p{(\linewidth - 16\tabcolsep) * \real{0.1071}}
  >{\raggedright\arraybackslash}p{(\linewidth - 16\tabcolsep) * \real{0.1071}}
  >{\raggedright\arraybackslash}p{(\linewidth - 16\tabcolsep) * \real{0.1071}}
  >{\raggedright\arraybackslash}p{(\linewidth - 16\tabcolsep) * \real{0.1071}}
  >{\raggedright\arraybackslash}p{(\linewidth - 16\tabcolsep) * \real{0.1071}}@{}}
\toprule\noalign{}
\begin{minipage}[b]{\linewidth}\raggedleft
CC
\end{minipage} & \begin{minipage}[b]{\linewidth}\raggedright
Function
\end{minipage} & \begin{minipage}[b]{\linewidth}\raggedright
Drum 1
\end{minipage} & \begin{minipage}[b]{\linewidth}\raggedright
Drum 2
\end{minipage} & \begin{minipage}[b]{\linewidth}\raggedright
Drum 3
\end{minipage} & \begin{minipage}[b]{\linewidth}\raggedright
Snare
\end{minipage} & \begin{minipage}[b]{\linewidth}\raggedright
Cymbal
\end{minipage} & \begin{minipage}[b]{\linewidth}\raggedright
HH A
\end{minipage} & \begin{minipage}[b]{\linewidth}\raggedright
HH B
\end{minipage} \\
\midrule\noalign{}
\endhead
\bottomrule\noalign{}
\endlastfoot
0 & Bank change & Yes & Yes & Yes & Yes & Yes & Yes & Yes \\
1 & Morph/mod wheel & Yes & Yes & Yes & Yes & Yes & Yes & Yes \\
7 & Voice volume & Yes & Yes & Yes & Yes & Yes & Yes & Yes \\
9 & Osc 1 waveform & Yes & Yes & Yes & Yes & Yes & Yes & Yes \\
10 & Pan & Yes & Yes & Yes & Yes & Yes & Yes & Yes \\
12 & Decimation & Yes & Yes & Yes & Yes & Yes & Yes & Yes \\
13 & Distortion & Yes & Yes & Yes & Yes & Yes & Yes & Yes \\
14 & Osc 1 coarse tune & Yes & Yes & Yes & Yes & Yes & Yes & Yes \\
16 & Filter frequency & Yes & Yes & Yes & Yes & Yes & Yes & Yes \\
17 & Filter resonance & Yes & Yes & Yes & Yes & Yes & Yes & Yes \\
18 & Filter drive & Yes & Yes & Yes & Yes & Yes & Yes & Yes \\
19 & Filter type & Yes & Yes & Yes & Yes & Yes & Yes & Yes \\
20 & Snare mix & -- & -- & -- & Yes & -- & -- & -- \\
21 & Snare noise frequency & -- & -- & -- & Yes & -- & -- & -- \\
22 & Mix/mod select & Yes & Yes & Yes & -- & -- & -- & -- \\
23 & Velocity volume mod on/off & Yes & Yes & Yes & Yes & Yes & Yes &
Yes \\
24 & Velocity mod amount & Yes & Yes & Yes & Yes & Yes & Yes & Yes \\
25 & Velocity mod destination & Yes & Yes & Yes & Yes & Yes & Yes &
Yes \\
26 & Transient volume & Yes & Yes & Yes & Yes & Yes & Yes & Yes \\
27 & Transient waveform & Yes & Yes & Yes & Yes & Yes & Yes & Yes \\
28 & Transient frequency & Yes & Yes & Yes & Yes & Yes & Yes & Yes \\
46 & Osc 1 fine tune & Yes & Yes & Yes & Yes & Yes & Yes & Yes \\
70 & Volume envelope repeat & -- & -- & -- & Yes & Yes & -- & -- \\
72 & Open hihat decay & -- & -- & -- & -- & -- & Yes & Yes \\
73 & Volume envelope attack & Yes & Yes & Yes & Yes & Yes & Yes & Yes \\
74 & Volume envelope slope & Yes & Yes & Yes & Yes & Yes & Yes & Yes \\
75 & Volume envelope decay & Yes & Yes & Yes & Yes & Yes & Yes & Yes \\
76 & LFO frequency & Yes & Yes & Yes & Yes & Yes & Yes & Yes \\
77 & LFO amount & Yes & Yes & Yes & Yes & Yes & Yes & Yes \\
78 & LFO phase offset & Yes & Yes & Yes & Yes & Yes & Yes & Yes \\
79 & LFO waveform & Yes & Yes & Yes & Yes & Yes & Yes & Yes \\
80 & LFO target voice selector & Yes & Yes & Yes & Yes & Yes & Yes &
Yes \\
81 & LFO target selector & Yes & Yes & Yes & Yes & Yes & Yes & Yes \\
82 & LFO retrigger & Yes & Yes & Yes & Yes & Yes & Yes & Yes \\
83 & LFO sync & Yes & Yes & Yes & Yes & Yes & Yes & Yes \\
84 & Mod envelope decay & Yes & Yes & Yes & Yes & -- & -- & -- \\
85 & Envelope mod amount & Yes & Yes & Yes & Yes & -- & -- & -- \\
86 & Mod envelope slope & Yes & Yes & Yes & Yes & -- & -- & -- \\
89 & Audio output routing & Yes & Yes & Yes & Yes & Yes & Yes & Yes \\
90 & Envelope position reset & Yes & Yes & Yes & Yes & Yes & Yes &
Yes \\
102 & Mod oscillator waveform & Yes & Yes & Yes & -- & -- & -- & -- \\
103 & FM/mod oscillator coarse tune & Yes & Yes & Yes & -- & -- & -- &
-- \\
104 & FM amount & Yes & Yes & Yes & -- & -- & -- & -- \\
105 & Mod oscillator waveform 1 & -- & -- & -- & -- & Yes & Yes & Yes \\
106 & Mod oscillator gain 1 & -- & -- & -- & -- & Yes & Yes & Yes \\
107 & Mod oscillator coarse tune 1 & -- & -- & -- & -- & Yes & Yes &
Yes \\
108 & Mod oscillator waveform 2 & -- & -- & -- & -- & Yes & Yes & Yes \\
109 & Mod oscillator gain 2 & -- & -- & -- & -- & Yes & Yes & Yes \\
110 & Mod oscillator coarse tune 2 & -- & -- & -- & -- & Yes & Yes &
Yes \\
120 & Track mute & Yes & Yes & Yes & Yes & Yes & Yes & Yes \\
\end{longtable}
}

\subsubsection{Global MIDI CC Table}\label{global-midi-cc-table}

This table applies only on the configured Global MIDI channel.

Notes on the CC and NRPN tables below:

\begin{itemize}
\item
  \textbf{Sending an NRPN:} send CC99 (NRPN number MSB), then CC98 (NRPN
  number LSB), then CC6 (value). CC6 is ignored until an NRPN number has
  been selected.
\item
  \textbf{Performance NRPNs:}

  {\def\LTcaptype{none} % do not increment counter
  \begin{longtable}[]{@{}
    >{\raggedleft\arraybackslash}p{(\linewidth - 4\tabcolsep) * \real{0.4000}}
    >{\raggedright\arraybackslash}p{(\linewidth - 4\tabcolsep) * \real{0.3000}}
    >{\raggedright\arraybackslash}p{(\linewidth - 4\tabcolsep) * \real{0.3000}}@{}}
  \toprule\noalign{}
  \begin{minipage}[b]{\linewidth}\raggedleft
  NRPN
  \end{minipage} & \begin{minipage}[b]{\linewidth}\raggedright
  Function
  \end{minipage} & \begin{minipage}[b]{\linewidth}\raggedright
  Range
  \end{minipage} \\
  \midrule\noalign{}
  \endhead
  \bottomrule\noalign{}
  \endlastfoot
  93 & Roll rate (Performance menu \param{rol}) & 0--15 \\
  94 & Roll note (Performance menu \param{nte}) & 0--127 \\
  95 & Roll velocity (Performance menu \param{vel}) & 0--127 \\
  96 & Filter frequency, voice 1 (¢ receive filter entries only) &
  0--127 \\
  \end{longtable}
  }

  NRPN 93--95 update the Performance menu display. They are not recorded
  as automation.
\item
  \textbf{CC38 swap option:} many devices, including Elektron gear, send
  CC38 (data entry LSB) after every NRPN value. On the global channel
  CC38 is voice 1 filter frequency, so every NRPN from such a device
  also closes voice 1's filter. Selecting one of the ¢ entries of the
  Receive MIDI Filter (\param{rxf}) makes the LXR ignore CC38 on the
  global channel and moves voice 1 filter frequency to NRPN 96. With
  standard entries, CC38 keeps its normal function and NRPN 96 is
  ignored. See \emph{Receive MIDI Filter} in the Global Settings Menu
  section.
\item
  \textbf{NRPN deselect (RPN null):} CC101 = 127 or CC100 = 127 sent
  directly after an NRPN value is treated as the standard NRPN deselect.
  It does not change voice pan, and a following CC6 is ignored until a
  new NRPN number is selected. CC100 and CC101 sent at any other time
  control voice pan as listed.
\end{itemize}
